\Section{Related Work and Conclusion}
\label{sec:conclusion}

\SubSection{Related work}
\label{sec:related}

Higher-order programs with side effects have been a long-standing concern in program verification.
Our approach to behavioral predicates descends from the relational tradition of VDM~\cite{jones1990vdm}, Z~\cite{spivey1992z, grieskamp2000executablez}, and the B method~\cite{abrial1996bbook}, in which operations are specified by relations between pre- and post-states; state labels generalize this to intermediate snapshots and make the relational flavor composable.

Among contemporary verifiers for imperative code with higher-order functions, \verus~\cite{verus} offers function-value specifications via \emph{spec closures} and phantom ghost state; the approach is state-passing, with state threaded as an extra argument to specifications rather than quantified over directly.
Our treatment is more directly relational: state labels denote memory snapshots and are bound by quantifiers in the surface syntax, avoiding the encoding cost of an explicit state argument and matching the style in which Move developers already write pre-/postcondition specifications.
Also for Rust, Wolff et al.~\cite{RustClosures21} verify closures in the Prusti verifier via \emph{call descriptions}, which relate a closure call's arguments and results to its captured state within separation logic; our behavioral predicates play a similar role, but are relational over labeled memory snapshots, with Move's alias-free references and resource-separated memory removing the need for separation-logic machinery.

\fstar~\cite{fstar2016} likewise verifies higher-order, effectful code, internalizing the WP calculus in its type system via Dijkstra monads~\cite{dm4all}.
Both systems support pre-/postcondition specifications, quantifiers, behavioral abstraction over function values, and structured proof guidance --- in Move via \texttt{proof} blocks, \texttt{lemma} declarations, and \texttt{calc}, \texttt{split}, \texttt{apply} statements.
The treatment of memory differs: \fstar~optionally supports separation logic (via the Steel and Pulse frameworks) to reason about heap separation within the logic, whereas Move uses traditional frame conditions via modifies declarations, and its resource and reference model separates memory syntactically at compile time --- separation need not be reasoned about by the SMT solver.

Dafny~\cite{DAFNY} is a close cousin of \MVP: it has a similar specification model of pre-/postconditions, higher-order functions, and behavioral predicates.
However, Dafny does not allow function values to modify state.
This would be a severe restriction for Move, since it is a common usage pattern.
Moreover, Move's resource-separated global memory and alias-free mutable reference model is not available in Dafny, as is also the case in \fstar.
While both languages support state labels, Move state labels as described here can be used in abstract specifications via universal and existential quantification.

Lean~4~\cite{lean4} is a dependently-typed interactive theorem prover~\cite{Isabelle, COQ} supporting higher-order functions, with tactic-driven, kernel-checked verification.
Unlike \MVP, Lean does not work directly on a programming language but requires the user to encode it in its meta logic, typically modeling stateful systems via monads.
\MVP, in contrast, integrates directly with Move and aims at automated SMT-based verification, with optional proof scripts.

In the smart-contract setting, dynamic dispatch is a central challenge for \solidity verifiers.
\certora~\cite{certora} handles dispatch on Solidity by enumerating the possible callees of a virtual call (derived from the contracts in scope) and generating a case-split proof obligation per callee; summaries can be supplied by the user to abstract external calls.
Other SMT-based Solidity tools~\cite{smtalt18,solcverify,DBLP:conf/esop/HajduJ20} handle dynamic dispatch through similar case-splitting, with varying degrees of abstraction.
\MVP is more general: a call site's dispatch is explicit in the closure's datatype, and our encoding lets the prover automatically fall back from case-split (when the concrete function is known) to behavioral-predicate reasoning (when only the function type is known), without the user having to structure the specification differently.

\SubSection{Future Work}

The current encoding enumerates all closure variants reaching a given call site; scaling to programs where the variant set is large or open (e.g.\ dynamically loaded modules) remains future work.
State labels are currently individually bound; extending quantification to \emph{sequences} of states ($\forall S_0..S_n$) would allow direct contracts for fold-like iteration, where today the higher-order function is verified once against a summary spec (Sec.~\ref{sec:find}) --- the flat existential encoding of Sec.~\ref{sec:enc-frame} already handles chains of fixed length and would extend naturally.
Proof automation for non-linear arithmetic arising in pricing formulas and similar domains currently relies on user-supplied proof hints; integrating portfolio arithmetic decision procedures is a natural direction.
The spec inference pipeline produces behavioral-predicate-based specifications from arbitrary code; extending it to synthesize loop invariants involving function values automatically would reduce the remaining manual annotation burden.

\SubSection{Conclusion}

We extended \MVP to Move 2's higher-order functions and dynamic dispatch.
The extension rests on two language additions --- behavioral predicates and state labels --- and on a Boogie encoding that discriminates closure invocations on the function-value variant reaching the call site.
The design preserves \MVP's production-oriented cost model: routine, predictable, counterexample-driven verification, also in the presence of dynamic dispatch.
We further extended \MVP's specification inference to closures and dynamic dispatch, turning compliance of a function value with a behavioral invariant into an automated verification problem.

%%% Local Variables:
%%% mode: latex
%%% TeX-master: "main"
%%% End:
